\begin{figure}[ht]
\centering
\begin{tikzpicture}[x=1.0cm, y=1.0cm, font=\small,
  sector/.style={rectangle, rounded corners=2pt, draw, thick,
    fill=engblue!12, minimum width=22mm, minimum height=10mm,
    align=center, inner sep=2pt}]

\node[sector] (agr) at (0, 0) {Agriculture};
\node[sector] (man) at (5, 0) {Manufacturing};
\node[sector] (ene) at (2.5, -3) {Energy};

% Inter-industry flows (a_ij = input from i needed per unit output of j)
\draw[->, thick, >=stealth] (agr) to[bend left=12]
  node[midway, above, font=\scriptsize] {$0.1$} (man);
\draw[->, thick, >=stealth] (man) to[bend left=12]
  node[midway, below, font=\scriptsize] {$0.2$} (agr);
\draw[->, thick, >=stealth] (ene) to[bend left=10]
  node[midway, left=1pt, font=\scriptsize] {$0.3$} (agr);
\draw[->, thick, >=stealth] (ene) to[bend right=10]
  node[midway, right=1pt, font=\scriptsize] {$0.2$} (man);
\draw[->, thick, >=stealth] (man) to[bend left=10]
  node[midway, right=1pt, font=\scriptsize] {$0.1$} (ene);
\draw[->, thick, engred, >=stealth] (agr) edge[loop above]
  node[font=\scriptsize] {$0.2$} (agr);

% Final demand arrows out
\draw[->, very thick, enggray, >=stealth] (agr) -- ++(-2.0, 0.0)
  node[left, font=\scriptsize] {demand $d_1$};
\draw[->, very thick, enggray, >=stealth] (man) -- ++(2.0, 0.0)
  node[right, font=\scriptsize] {demand $d_2$};
\draw[->, very thick, enggray, >=stealth] (ene) -- ++(0.0, -1.6)
  node[below, font=\scriptsize] {demand $d_3$};

\node[anchor=north, font=\scriptsize, text width=10.5cm, align=center]
  at (2.5, -5.2) {
  A three-sector Leontief economy. An arrow labelled $a_{ij}$ from
  sector $i$ to sector $j$ means producing one unit of sector $j$'s
  output consumes $a_{ij}$ units of sector $i$'s output. The
  technical-coefficient matrix $\mat{A}$ collects these labels; the
  total output $\vect{x}$ that meets a final demand $\vect{d}$ solves
  $(\mat{I} - \mat{A})\vect{x} = \vect{d}$.};

\end{tikzpicture}
\caption[A Leontief input-output economy as a linear
system]{A three-sector input-output economy. Each sector's output is
consumed partly by the other sectors as intermediate input and partly
by final demand. The technical-coefficient matrix $\mat{A}$ records
the inter-industry requirements, and the gross output that meets a
final demand $\vect{d}$ solves the Leontief system $(\mat{I} -
\mat{A})\vect{x} = \vect{d}$. The economy is productive (the system
has a non-negative solution for every non-negative demand) exactly
when the spectral radius of $\mat{A}$ is below one, which is the
convergence condition for the Neumann series $(\mat{I} - \mat{A})^{-1}
= \sum_{k \geq 0} \mat{A}^{k}$.}
\label{fig:vol02:ch09:leontief}
\end{figure}
